{\large\textbf{Champagne! (10 points)}}

\textbf{Warning:} Excessive alcohol consumption is harmful to health and
drinking alcohol below legal age is prohibited.

Champagne is a French sparkling wine. Fermentation of sugars produces carbon
dioxide ($\mathrm{CO_2}$) in the bottle. The molar concentration of
$\mathrm{CO_2}$ in the liquid phase $c_\ell$ and the partial pressure
$P_{\mathrm{CO_2}}$ in the gas phase are related by
$c_\ell = k_\mathrm{H}P_{\mathrm{CO_2}}$, known as Henry's law and where
$k_\mathrm{H}$ is called Henry's constant.

\textbf{Data}
\begin{itemize}
\item Surface tension of champagne
$\sigma = 47 \times 10^{-3}\,\mathrm{J\cdot m^{-2}}$
\item Density of the liquid
$\rho_\ell = 1.0 \times 10^3\,\mathrm{kg\cdot m^{-3}}$
\item Henry's constant at $T_0 = 20\,^\circ\mathrm{C}$,
$k_\mathrm{H}(20\,^\circ\mathrm{C}) = 3.3 \times 10^{-4}\,\mathrm{mol\cdot m^{-3}\cdot Pa^{-1}}$
\item Henry's constant at $T_0 = 6\,^\circ\mathrm{C}$,
$k_\mathrm{H}(6\,^\circ\mathrm{C}) = 5.4 \times 10^{-4}\,\mathrm{mol\cdot m^{-3}\cdot Pa^{-1}}$
\item Atmospheric pressure
$P_0 = 1\,\mathrm{bar} = 1.0 \times 10^5\,\mathrm{Pa}$
\item Gases are ideal with an adiabatic coefficient $\gamma = 1.3$
\end{itemize}

\begin{center}\includegraphics[width=0.09\linewidth]{figures/IPhO_2025_Q3_fig1.jpg}\end{center}

\begin{center}\textbf{Fig.~1}. A glass filled with champagne.\end{center}

\textbf{Part A. Nucleation, growth and rise of bubbles}

Immediately after opening a bottle of champagne at temperature
$T_0 = 20\,^\circ\mathrm{C}$ , we fill a glass. The pressure in the liquid is
$P_0$ and its temperature stays constant at $T_0$. The concentration $c_\ell$
of dissolved $\mathrm{CO_2}$ exceeds the equilibrium concentration and we study
the nucleation of a $\mathrm{CO_2}$ bubble. We note $a$ its radius and
$P_\mathrm{b}$ its inner pressure.

\textbf{A.1}\quad Express the pressure $P_\mathrm{b}$ in terms of $P_0$, $a$
and $\sigma$. \hfill 0.2pt

In the liquid, the concentration of dissolved $\mathrm{CO_2}$ depends on the
distance to the bubble. At long distance we recover the value $c_\ell$ and we
note $c_\mathrm{b}$ the concentration close to the bubble surface. According to
Henry's law, $c_\mathrm{b} = k_\mathrm{H}P_\mathrm{b}$. We furthermore assume
in all the problem that bubbles contain only $\mathrm{CO_2}$.

Since $c_\ell \neq c_\mathrm{b}$, $\mathrm{CO_2}$ molecules diffuse from areas
of high to low concentration. We assume also that any molecule from the liquid
phase reaching the bubble surface is transferred to the vapour.

\textbf{A.2}\quad Express the critical radius $a_\mathrm{c}$ above which a
bubble is expected to grow in terms of $P_0,\sigma,c_\ell$ and $c_0$ where
$c_0 = k_\mathrm{H}P_0$. Calculate numerically $a_\mathrm{c}$ for
$c_\ell = 4c_0$. \hfill 0.5pt

In practice, bubbles mainly grow from pre-existing gas cavities. Consider then
a bubble with initial radius $a_0 \approx 40\,\upmu\mathrm{m}$. The number of
moles of $\mathrm{CO_2}$ transferred at the bubble's surface per unit area and
time is noted $j$. Two models are possible for $j$.
\begin{itemize}
\item model (1) $j = \dfrac{D}{a}(c_\ell - c_\mathrm{b})$ where $D$ is the
diffusion coefficient of $\mathrm{CO_2}$ in the liquid.
\item model (2) $j = K(c_\ell - c_\mathrm{b})$ where $K$ is a constant here.
\end{itemize}

Experimentally, the bubble radius $a(t)$ is found to depend on time as shown in
\textbf{Fig.~2}. Here $c_\ell \approx 4c_0$, and since bubbles are large enough
to be visible, the excess pressure due to surface tension can be neglected and
$P_\mathrm{b} \approx P_0$.

\textbf{A.3}\quad Express the number of $\mathrm{CO_2}$ moles in the bubble
$n_\mathrm{c}$ in terms of $a, P_0, T_0$ and ideal gas constant $R$. Find
$a(t)$ for both models. Indicate which model explains the experimental results
in \textbf{Fig.~2}. Depending on your answer, calculate numerically $K$ or
$D$. \hfill 1.2pt

\begin{center}\includegraphics[width=0.45\linewidth]{figures/IPhO_2025_Q3_fig2.png}\end{center}

\begin{center}\textbf{Fig.~2.} Time evolution of $\mathrm{CO_2}$ bubble radius
in a glass of champagne (\textit{adapted from [1]}).\end{center}

Eventually bubbles detach from the bottom of the glass and continue to grow
while rising. \textbf{Fig.~3}. shows a train of bubbles. The bubbles of the
train have the same initial radius and are emitted at a constant frequency
$f_\mathrm{b} = 20\,\mathrm{Hz}$.

\begin{center}\includegraphics[width=\linewidth]{figures/IPhO_2025_Q3_fig3.png}\end{center}

\begin{center}\textbf{Fig.~3.} A train of bubbles. The photo is rotated
horizontally for the page layout \textit{(adapted from [1]).}\end{center}

For the range of velocities studied here, the drag force $F$ on a bubble of
radius $a$ moving at velocity $v$ in a liquid of dynamic viscosity $\eta$ is
given by Stokes' law $F = 6\pi\eta av$. Measurements show that at any moment in
time, the bubble can be assumed to be travelling at its terminal velocity.

\textbf{A.4}\quad Give the expression of the main forces exerted on a
vertically rising bubble. Obtain the expression of $v(a)$. Give a numerical
estimate of $\eta$ using $\rho_\ell$, $g_0$ and quantities measured on
\textbf{Fig.~3.} \hfill 0.8pt

The quasi-stationary growth of bubbles with rate
$q_a = \frac{\mathrm{d}a}{\mathrm{d}t}$ still applies during bubble rise.

\textbf{A.5}\quad Express the radius $a_{H_\ell}$ of a bubble reaching the free
surface in terms of height travelled $H_\ell$, growth rate
$q_a = \frac{\mathrm{d}a}{\mathrm{d}t}$, and any constants you may need.
Assume $a_{H_\ell} \gg a_0$ and $q_a$ constant, and give the numerical value of
$a_{H_\ell}$ with $H_\ell = 10\,\mathrm{cm}$ and $q_a$ corresponding to
\textbf{Fig.~2}. \hfill 0.5pt

There are $N_\mathrm{b}$ nucleation sites of bubbles. Assume that the bubbles
are nucleated at a constant frequency $f_\mathrm{b}$ at the bottom of a glass of
champagne (height $H_\ell$ for a volume $V_\ell$), with $a_0$ still negligible.
Neglect diffusion of $\mathrm{CO_2}$ at the free surface.

\textbf{A.6}\quad Write the differential equation for $c_\ell(t)$. Obtain from
this equation the characteristic time $\tau$ for the decay of the concentration
of dissolved $\mathrm{CO_2}$ in the liquid. \hfill 1.1pt

\textbf{Part B. Acoustic emission of a bursting bubble}

Small bubbles are nearly spherical as they reach the free surface. Once the
liquid film separating the bubble from the air thins out sufficiently, a
circular hole of radius $r$ forms in the film and, driven by surface tension,
opens very quickly (\textbf{Fig.~4.} left). The hole opens at constant speed
$v_\mathrm{f}$ (\textbf{Fig.~4.} right). The film outside the rim remains
still, with constant thickness $h$.

\begin{center}\includegraphics[width=0.79\linewidth]{figures/IPhO_2025_Q3_fig4.png}\end{center}

\textbf{Fig.~4}. \textit{(Left)} $(\alpha)$ Bubble at the surface: (1) liquid,
(2) air at pressure $P_0$ and (3), $\mathrm{CO_2}$ at pressure $P_\mathrm{b}$,
$(\beta)$ and $(\gamma)$ retraction of the liquid film, where the rim is in dark
blue, $(\delta)$ bubble collapse. \textit{(Right)} Retraction of the liquid film
at time $t$. Top: sketch of the pierced film seen from above. Bottom:
cross-section of the rim and the retracting film. During $\mathrm{d}t$ the rim
accumulates nearby liquid (dotted).

Due to dissipative processes, only half of the difference of the surface energy
between $t$ and $t+\mathrm{d}t$ of the rim and the accumulated liquid is
transformed into kinetic energy. We further assume that the variation of the
surface of the rim is negligible compared to that of the film.

\textbf{B.1}\quad Express $v_\mathrm{f}$ in terms of $\rho_\ell,\sigma$ and
$h$. \hfill 1.1pt

\begin{center}\includegraphics[width=0.47\linewidth]{figures/IPhO_2025_Q3_fig5.png}\end{center}

\begin{center}\textbf{Fig.~5}. \textit{(Left)} a Helmholtz resonator.
\textit{(Right)} a bubble as an oscillator.\end{center}

When the film bursts, it releases internal pressure and emits a sound. We model
this acoustic emission by a Helmholtz resonator: a cavity open to the
atmosphere at $P_0$ through a bottleneck aperture of area $S$
(\textbf{Fig.~5}. left). In the neck, a mass $m_\mathrm{p}$ makes small
amplitude position oscillations due to the pressure forces it experiences as
the gas in the cavity expands or compresses adiabatically. The gravity force on
$m_\mathrm{p}$ is negligible compared to pressure forces. Let $V_0$ be the
volume of gas under the mass $m_\mathrm{p}$ for $P = P_0$ as $z = 0$.

\textbf{B.2}\quad Express the frequency of oscillation $f_0$ of $m_\mathrm{p}$.
Hint: for $\varepsilon \ll 1$, $(1 + \varepsilon)^\alpha \approx 1 +
\alpha\varepsilon$. \hfill 1.1pt

The Helmholtz model may be used for a bubble of radius $a$. $V_0$ is the volume
of the closed bubble. From litterature, the mass of the equivalent of the
piston is $m_p = 8\rho_g r^3/3$ where $r$ is the radius of the circular
aperture and $\rho_g = 1.8\,\mathrm{kg\cdot m^{-3}}$ is the density of the gas
(\textbf{Fig.~5}. right). During the bursting process, $r$ goes from 0 to
$r_\mathrm{c}$, given by
$r_\mathrm{c} = \dfrac{2}{\sqrt{3}}a^2\sqrt{\dfrac{\rho_\ell g_0}{\sigma}}$. At
the same time, the frequency of emitted sound increases until a maximum value
of $40\,\mathrm{kHz}$ and the bursting time is
$t_b = 3 \times 10^{-2}\,\mathrm{ms}$.

\textbf{B.3}\quad Find the radius $a$ and the thickness $h$ of the champagne
film separating the bubble from the atmosphere. \hfill 1.1pt

\textbf{Part C. Popping champagne}

In a bottle, the total quantity of $\mathrm{CO_2}$ is
$n_\mathrm{T} = 0.2\,\mathrm{mol}$, either dissolved in the volume
$V_\mathrm{L} = 750\,\mathrm{mL}$ of liquid champagne, or as a gas in the
volume $V_\mathrm{G} = 25\,\mathrm{mL}$ under the cork (\textbf{Fig.~6.} left).
$V_\mathrm{G}$ contains only $\mathrm{CO_2}$. The equilibrium between both
$\mathrm{CO_2}$ phases follows Henry's Law. We suppose that the fast gaseous
$\mathrm{CO_2}$ expansion when the bottle is opened, is adiabatic and
reversible. Ambient temperature $T_0$ and pressure $P_0 = 1\,\mathrm{bar}$ are
constant.

\begin{center}\includegraphics[width=0.61\linewidth]{figures/IPhO_2025_Q3_fig6.png}\end{center}

\textbf{Fig.~6.} \textit{Left:} traditional bottleneck: (1) surrounding air,
(2) cork stopper, (3) headspace, (4) liquid champagne. \textit{Right}: Two
phenomena observed while opening the bottle at two different temperatures
(adapted from \textit{[2]}).

\textbf{C.1}\quad Give the numerical value of the pressure $P_\mathrm{i}$ of
gaseous $\mathrm{CO_2}$ in the bottle for $T_0 = 6\,^\circ\mathrm{C}$ and
$T_0 = 20\,^\circ\mathrm{C}$. \hfill 0.4pt

Another step of champagne production (not described here) leads to the
following values of $P_i$ that we will use for the next questions:
$P_\mathrm{i} = 4.69\,\mathrm{bar}$ at $T_0 = 6\,^\circ\mathrm{C}$ and
$P_\mathrm{i} = 7.45\,\mathrm{bar}$ at $T_0 = 20\,^\circ\mathrm{C}$.

During bottle opening, two different phenomena can be observed, depending on
$T_0$ (Fig.~6. right).
\begin{itemize}
\item either a blue fog appears, due to the formation of solid $\mathrm{CO_2}$
crystals (but water condensation is inhibited);
\item or a grey-white fog appears, due to water vapor condensation in the air
surrounding the bottleneck. In this latter case, there is no formation of
$\mathrm{CO_2}$ solid crystals.
\end{itemize}

The saturated vapor pressure $P_\mathrm{sat}^{\mathrm{CO_2}}$ for the
$\mathrm{CO_2}$ solid/gas transition follows :
$\log_{10}\left(\dfrac{P_\mathrm{sat}^{\mathrm{CO_2}}}{P_0}\right) = A -
\dfrac{B}{T + C}$ with $T$ in $\mathrm{K}$, $A = 6.81$,
$B = 1.30 \times 10^3\,\mathrm{K}$ and $C = -3.49\,\mathrm{K}$.

\textbf{C.2}\quad Give the numerical value $T_\mathrm{f}$ of the
$\mathrm{CO_2}$ gas at the end of the expansion, after opening a bottle, if
$T_0 = 6\,^\circ\mathrm{C}$ and if $T_0 = 20\,^\circ\mathrm{C}$, if no phase
transition occured. Choose which statements are true (several statements
possible):\\
\hspace*{1.5em}1.\ At $T_0 = 6\,^\circ\mathrm{C}$ a grey-white fog appears
while opening the bottle.\\
\hspace*{1.5em}2.\ At $T_0 = 6\,^\circ\mathrm{C}$ a blue fog appears while
opening the bottle.\\
\hspace*{1.5em}3.\ At $T_0 = 20\,^\circ\mathrm{C}$ a grey-white fog appears
while opening the bottle.\\
\hspace*{1.5em}4.\ At $T_0 = 20\,^\circ\mathrm{C}$ a blue fog appears while
opening the bottle. \hfill 0.7pt

During bottle opening, the cork stopper pops out. We now determine the maximum
height $H_\mathrm{c}$ it reaches. Assume that the friction force $F$ due to the
bottleneck on the cork stopper is $F = \alpha A$ where $A$ is the area of
contact and $\alpha$ is a constant to determine. Initially, the pressure force
slightly overcomes the friction force. The cork's mass is
$m = 10\,\mathrm{g}$, its diameter $d = 1.8\,\mathrm{cm}$ and the length of
the cylindrical part initially stuck in the bottleneck is
$\ell_0 = 2.5\,\mathrm{cm}$. Once the cork has left the bottleneck, you can
neglect the net pressure force.

\textbf{C.3}\quad Give the numerical value of $H_\mathrm{c}$ if the external
temperature is $T_0 = 6\,^\circ\mathrm{C}$. \hfill 1.3pt

[1] Liger-Belair \textit{et al,} Am. J. Enol. Vitic., Vol. 50, No. 3 (1999).

[2] Liger-Belair \textit{et al.}, \textit{Sc. Reports} \textbf{7}, 10938
(2017).
